\documentclass{beamer}

% Theme choice
\usetheme{Madrid} % You can change to e.g., Warsaw, Berlin, CambridgeUS, etc.

% Encoding and font
\usepackage[utf8]{inputenc}
\usepackage[T1]{fontenc}

% Graphics and tables
\usepackage{graphicx}
\usepackage{booktabs}

% Code listings
\usepackage{listings}
\lstset{
basicstyle=\ttfamily\small,
keywordstyle=\color{blue},
commentstyle=\color{gray},
stringstyle=\color{red},
breaklines=true,
frame=single
}

% Math packages
\usepackage{amsmath}
\usepackage{amssymb}

% Colors
\usepackage{xcolor}

% TikZ and PGFPlots
\usepackage{tikz}
\usepackage{pgfplots}
\pgfplotsset{compat=1.18}
\usetikzlibrary{positioning}

% Hyperlinks
\usepackage{hyperref}

% Title information
\title{Week 11: Software Design Patterns}
\author{Your Name}
\institute{Your Institution}
\date{\today}

\begin{document}

\frame{\titlepage}

\begin{frame}[fragile]
    \frametitle{Introduction to Software Design Patterns}
    \begin{block}{Overview}
        Software design patterns are general reusable solutions to common problems in software design. They provide a standardized method to address issues that arise during development, making it easier to communicate ideas among developers and to recognize proven techniques.
    \end{block}
\end{frame}

\begin{frame}[fragile]
    \frametitle{Importance of Design Patterns}
    \begin{enumerate}
        \item \textbf{Enhanced Communication:}
        Patterns establish a common language among developers, allowing for clearer discussions of design. For example, using terms like "Singleton" or "Observer" communicates specific design strategies without lengthy explanations.
        
        \item \textbf{Improved Code Maintainability:}
        Design patterns make code more understandable and manageable, leading to more robust and maintainable codebases.

        \item \textbf{Increased Reusability:}
        By applying design patterns, developers can create components that can be reused across different projects, ensuring consistency in solutions.

        \item \textbf{Facilitates Design Flexibility:}
        Patterns encourage the use of interfaces and abstraction, enabling systems to be designed for change over time.

        \item \textbf{Encourages Best Practices:}
        Design patterns encapsulate best practices, guiding developers in making solid architectural decisions.
    \end{enumerate}
\end{frame}

\begin{frame}[fragile]
    \frametitle{Key Points to Emphasize}
    \begin{itemize}
        \item Design patterns are templates that can be adapted to specific contexts.
        \item Familiarity with design patterns can enhance a developer’s speed in producing high-quality software.
        \item Patterns can be categorized into:
        \begin{itemize}
            \item \textbf{Creational Patterns} (e.g., Singleton, Factory)
            \item \textbf{Structural Patterns} (e.g., Adapter, Composite)
            \item \textbf{Behavioral Patterns} (e.g., Observer, Strategy)
        \end{itemize}
    \end{itemize}
\end{frame}

\begin{frame}[fragile]
    \frametitle{Example: The Singleton Pattern}
    \begin{lstlisting}[language=Python]
class Singleton:
    _instance = None

    def __new__(cls):
        if cls._instance is None:
            cls._instance = super().__new__(cls)
        return cls._instance
    \end{lstlisting}
    \begin{block}{Explanation}
        In this example, the Singleton pattern ensures that a class has only one instance, providing a global point of access to it.
    \end{block}
\end{frame}

\begin{frame}[fragile]
    \frametitle{Conclusion}
    By understanding and leveraging software design patterns, developers can significantly improve their skills and the quality of their software solutions. Patterns serve as valuable tools in a developer’s arsenal, fostering a deeper understanding of design principles and best practices.
\end{frame}

\begin{frame}[fragile]
    \frametitle{What are Design Patterns?}
    \begin{block}{Definition of Design Patterns}
        Design patterns are reusable solutions to common problems that occur within a given context in software design. They provide a tried-and-true methodology for solving design issues and offer a shared language among developers to communicate complex designs effectively.
    \end{block}
\end{frame}

\begin{frame}[fragile]
    \frametitle{Purpose of Design Patterns}
    The primary purposes of design patterns in software development include:
    \begin{enumerate}
        \item \textbf{Standardization}: Provide standard terminology and concepts for collaboration.
        \item \textbf{Code Reusability}: Promote reuse across applications, reducing redundancy.
        \item \textbf{Best Practices}: Encapsulate proven methodologies for robust software solutions.
        \item \textbf{Maintainability}: Improve readability and ease future modifications.
        \item \textbf{Scalability}: Facilitate system extensions by promoting loose coupling.
    \end{enumerate}
\end{frame}

\begin{frame}[fragile]
    \frametitle{Examples of Design Patterns}
    Two common design patterns are:
    
    \begin{block}{Singleton Pattern}
        Ensures a class has only one instance and provides a global point of access. Useful for single configuration objects.
        \begin{lstlisting}[language=Python]
class Singleton:
    _instance = None

    def __new__(cls):
        if cls._instance is None:
            cls._instance = super(Singleton, cls).__new__(cls)
        return cls._instance
        \end{lstlisting}
    \end{block}

    \begin{block}{Observer Pattern}
        Defines a one-to-many dependency between objects, allowing one object to notify multiple others when its state changes.
        \begin{lstlisting}[language=Python]
class Subject:
    def __init__(self):
        self._observers = []

    def attach(self, observer):
        self._observers.append(observer)

    def notify(self):
        for observer in self._observers:
            observer.update()
        \end{lstlisting}
    \end{block}
\end{frame}

\begin{frame}[fragile]
    \frametitle{Benefits of Using Design Patterns - Introduction}
    Design patterns are established solutions to common software design problems. They enhance the effectiveness of software development by introducing best practices that lead to improved software qualities. This slide outlines the key benefits of using design patterns, namely code reusability, scalability, and maintainability.
\end{frame}

\begin{frame}[fragile]
    \frametitle{Benefits of Using Design Patterns - Improved Code Reusability}
    \begin{itemize}
        \item \textbf{Definition:} Reusability refers to the ability to use existing code in multiple applications, reducing redundancy and saving development time.
        \item \textbf{Explanation:} By utilizing design patterns, developers can create modular and standardized solutions that can be easily integrated into various programs.
        \item \textbf{Example:} The Factory Method pattern allows for the creation of objects without specifying the exact class of the object.
    \end{itemize}
    \begin{lstlisting}[language=Python]
class VehicleFactory:
    def create_vehicle(self, vehicle_type):
        if vehicle_type == "Car":
            return Car()
        elif vehicle_type == "Bike":
            return Bike()
    \end{lstlisting}
\end{frame}

\begin{frame}[fragile]
    \frametitle{Benefits of Using Design Patterns - Enhanced Scalability}
    \begin{itemize}
        \item \textbf{Definition:} Scalability is the ability of a software system to grow and manage increased demand.
        \item \textbf{Explanation:} Design patterns often encourage a decoupled architecture, allowing systems to scale more easily.
        \item \textbf{Example:} The Observer pattern lets your program communicate changes across different parts of the system without tight coupling.
    \end{itemize}
    \begin{lstlisting}[language=Python]
class Subject:
    def attach(self, observer):
        # Code to attach observer
    def notify_observers(self):
        # Code to notify all observers of changes
    \end{lstlisting}
\end{frame}

\begin{frame}[fragile]
    \frametitle{Benefits of Using Design Patterns - Improved Maintainability}
    \begin{itemize}
        \item \textbf{Definition:} Maintainability refers to the ease with which software can be updated or modified.
        \item \textbf{Explanation:} Design patterns provide a clear structure and guidelines that help future developers understand the code better.
        \item \textbf{Example:} The Model-View-Controller (MVC) pattern separates an application into three interconnected components.
    \end{itemize}
    \begin{block}{Application Architecture}
        - Model: Handles data and business logic \\
        - View: Responsible for UI and user interaction \\
        - Controller: Manages communication between Model and View
    \end{block}
\end{frame}

\begin{frame}[fragile]
    \frametitle{Benefits of Using Design Patterns - Key Points & Conclusion}
    \begin{itemize}
        \item \textbf{Reusability:} Save time and reduce errors by reusing code across projects.
        \item \textbf{Scalability:} Adapt quickly to growth by adding new features seamlessly.
        \item \textbf{Maintainability:} Simplify updates and debugging by utilizing structured patterns.
    \end{itemize}
    Utilizing design patterns fosters better software development practices by promoting code reusability, scalability, and maintainability.
\end{frame}

\begin{frame}[fragile]
    \frametitle{Benefits of Using Design Patterns - Next Steps}
    In the following slide, we will explore the Singleton Pattern, a widely-used design pattern that exemplifies these benefits.
\end{frame}

\begin{frame}[fragile]
    \frametitle{Singleton Pattern - Overview}
    \begin{block}{What is the Singleton Pattern?}
    The Singleton Pattern is a design pattern that restricts a class to a single instance while providing a global access point to that instance. This pattern is particularly useful when you need exactly one object to coordinate actions across the system.
    \end{block}
\end{frame}

\begin{frame}[fragile]
    \frametitle{Singleton Pattern - Key Characteristics}
    \begin{itemize}
        \item \textbf{Single Instance:} Only one instance of the class is created.
        \item \textbf{Global Access:} Provides a way to access that instance globally.
        \item \textbf{Lazy Initialization:} The instance is created only when needed, saving resources.
    \end{itemize}
\end{frame}

\begin{frame}[fragile]
    \frametitle{Singleton Pattern - Structure and Example}
    \begin{block}{Structure}
    \begin{enumerate}
        \item Private Constructor: Prevents external instantiation.
        \item Static Instance Variable: Holds the single instance.
        \item Public Static Method: Provides access to the instance.
    \end{enumerate}
    \end{block}

    \begin{lstlisting}[language=pseudocode]
class Singleton {
    private static Singleton instance;

    private Singleton() { }  // Private constructor

    public static Singleton getInstance() {
        if (instance == null) {
            instance = new Singleton(); // Lazy initialization
        }
        return instance;
    }
}
    \end{lstlisting}
\end{frame}

\begin{frame}[fragile]
    \frametitle{Singleton Pattern - Use Cases}
    \begin{itemize}
        \item \textbf{Logging:} Ensures a single point for logging messages.
        \item \textbf{Configuration Settings:} Provides consistent access to configuration data.
        \item \textbf{Thread Pool Management:} Manages threads efficiently with a single instance.
    \end{itemize}
\end{frame}

\begin{frame}[fragile]
    \frametitle{Singleton Pattern - Key Points}
    \begin{itemize}
        \item \textbf{Avoids Multiple Instances:} Prevents inconsistent behavior.
        \item \textbf{Eager vs Lazy Instantiation:} Understand the difference.
        \item \textbf{Thread Safety:} Caution is needed in multi-threaded applications.
    \end{itemize}
\end{frame}

\begin{frame}[fragile]
    \frametitle{Singleton Pattern - Summary and Next Steps}
    \begin{block}{Summary}
    The Singleton Pattern provides a controlled approach to instance management, leading to more efficient and manageable code.
    \end{block}

    \begin{block}{Next Steps}
    In the following slides, we will delve into implementing the Singleton Pattern in various popular programming languages.
    \end{block}
\end{frame}

\begin{frame}[fragile]
    \frametitle{Understanding the Singleton Pattern}
    The Singleton Pattern is a design pattern that restricts the instantiation of a class to a single instance. This is useful when exactly one object is needed to coordinate actions across the system.

    \begin{itemize}
        \item \textbf{Single Instance:} Ensures only one instance of a class is created.
        \item \textbf{Global Access Point:} Provides a global point of access to the instance.
        \item \textbf{Lazy Initialization:} Usually implemented with deferred instantiation of the instance.
    \end{itemize}
\end{frame}

\begin{frame}[fragile]
    \frametitle{Step-by-Step Implementation - Java}
    \begin{block}{Java}
    \begin{lstlisting}[language=Java]
public class Singleton {
    private static Singleton instance;

    private Singleton() {}

    public static Singleton getInstance() {
        if (instance == null) {
            instance = new Singleton();
        }
        return instance;
    }
}
    \end{lstlisting}
    \end{block}
    Example: Call \texttt{Singleton.getInstance()} to get the instance.
\end{frame}

\begin{frame}[fragile]
    \frametitle{Step-by-Step Implementation - Python, C\#, JavaScript}
    \begin{block}{Python}
    \begin{lstlisting}[language=Python]
class Singleton:
    _instance = None

    def __new__(cls):
        if cls._instance is None:
            cls._instance = super(Singleton, cls).__new__(cls)
        return cls._instance
    \end{lstlisting}
    \end{block}
    Example: Use \texttt{singleton = Singleton()} to access the singleton instance.

    \begin{block}{C\#}
    \begin{lstlisting}[language=CSharp]
public class Singleton {
    private static Singleton instance;
    private static readonly object lockObject = new object();

    private Singleton() {}

    public static Singleton Instance {
        get {
            lock (lockObject) {
                if (instance == null) {
                    instance = new Singleton();
                }
                return instance;
            }
        }
    }
}
    \end{lstlisting}
    \end{block}
   Example: Call \texttt{var singleton = Singleton.Instance} to retrieve the single instance.

   \begin{block}{JavaScript}
   \begin{lstlisting}[language=JavaScript]
const Singleton = (function() {
    let instance;

    function createInstance() {
        const object = new Object("I am the singleton instance");
        return object;
    }

    return {
        getInstance: function() {
            if (!instance) {
                instance = createInstance();
            }
            return instance;
        }
    };
})();
   \end{lstlisting}
   \end{block}
   Example: \texttt{const singleton = Singleton.getInstance();}
\end{frame}

\begin{frame}[fragile]
    \frametitle{Observer Pattern - Introduction}
    \begin{block}{Definition}
        The Observer Pattern is a behavioral design pattern that defines a one-to-many dependency between objects. When one object (the subject) changes its state, all its dependents (observers) are notified and updated automatically.
    \end{block}
    \begin{block}{Use Cases}
        - Useful when one part of a system must notify others about state changes.
        - Promotes loose coupling between components.
    \end{block}
\end{frame}

\begin{frame}[fragile]
    \frametitle{Observer Pattern - Structure}
    \begin{enumerate}
        \item \textbf{Subject}:
            \begin{itemize}
                \item Maintains a list of observers.
                \item Notifies observers about state changes.
            \end{itemize}
        \item \textbf{Observers}:
            \begin{itemize}
                \item Subscribe to the subject to receive updates.
                \item Implement an update interface.
            \end{itemize}
        \item \textbf{Concrete Subject}:
            \begin{itemize}
                \item Subclass of the subject that maintains state.
                \item Sends notifications to its observers.
            \end{itemize}
        \item \textbf{Concrete Observer}:
            \begin{itemize}
                \item Subclass of the observer that defines reactions to updates.
            \end{itemize}
    \end{enumerate}
\end{frame}

\begin{frame}[fragile]
    \frametitle{Observer Pattern - Example}
    \textbf{Scenario:} A weather station updating a display with temperature changes.
    \begin{block}{Weather Station Class}
    \begin{lstlisting}[language=Python]
class WeatherStation:
    def __init__(self):
        self._observers = []
        self._temperature = None

    def attach(self, observer):
        self._observers.append(observer)

    def detach(self, observer):
        self._observers.remove(observer)

    def notify(self):
        for observer in self._observers:
            observer.update(self._temperature)

    def set_temperature(self, temperature):
        self._temperature = temperature
        self.notify()
    \end{lstlisting}
    \end{block}
    \begin{block}{Temperature Display Class}
    \begin{lstlisting}[language=Python]
class TemperatureDisplay:
    def update(self, temperature):
        print(f"Temperature updated to: {temperature}°C")
    \end{lstlisting}
    \end{block}
    \textbf{Usage:}
    \begin{lstlisting}[language=Python]
weather_station = WeatherStation()
temp_display = TemperatureDisplay()

weather_station.attach(temp_display)
weather_station.set_temperature(25)  # Output: Temperature updated to: 25°C
    \end{lstlisting}
\end{frame}

\begin{frame}[fragile]
    \frametitle{Implementing the Observer Pattern - Overview}
    \begin{block}{Overview of the Observer Pattern}
        The Observer Pattern is a behavioral design pattern that establishes a one-to-many dependency between objects. 
        When one object (the subject) changes its state, all registered observers are notified and updated automatically.
        This pattern is widely used in scenarios where changes in one part of an application need to be reflected in another (e.g., in GUI toolkits, messaging systems).
    \end{block}
\end{frame}

\begin{frame}[fragile]
    \frametitle{Implementing the Observer Pattern - Key Components}
    \begin{block}{Key Components}
        \begin{itemize}
            \item \textbf{Subject}: The object that holds the state and notifies observers of any changes.
            \item \textbf{Observer}: An interface to be implemented by any class that wants to be updated when the subject changes.
            \item \textbf{ConcreteSubject}: A specific implementation of the subject that maintains the state and notifies observers.
            \item \textbf{ConcreteObserver}: An implementation of the observer that responds to changes in the subject.
        \end{itemize}
    \end{block}
\end{frame}

\begin{frame}[fragile]
    \frametitle{Implementing the Observer Pattern - Steps}
    \begin{block}{Steps to Implement the Observer Pattern}
        \begin{enumerate}
            \item \textbf{Define the Observer Interface}
            \begin{lstlisting}[language=Java]
    public interface Observer {
        void update(String message);
    }
            \end{lstlisting}

            \item \textbf{Create the Subject Class}
            \begin{lstlisting}[language=Java]
    import java.util.ArrayList;
    import java.util.List;

    public class Subject {
        private List<Observer> observers = new ArrayList<>();
        
        public void attach(Observer observer) {
            observers.add(observer);
        }
        
        public void detach(Observer observer) {
            observers.remove(observer);
        }
        
        public void notifyObservers(String message) {
            for (Observer observer : observers) {
                observer.update(message);
            }
        }
    }
            \end{lstlisting}

            \item \textbf{Implement the ConcreteSubject}
            \begin{lstlisting}[language=Java]
    public class WeatherStation extends Subject {
        private String weather;

        public void setWeather(String weather) {
            this.weather = weather;
            notifyObservers(weather);
        }
    }
            \end{lstlisting}
        \end{enumerate}
    \end{block}
\end{frame}

\begin{frame}[fragile]
    \frametitle{Implementing the Observer Pattern - Example Usage}
    \begin{block}{Putting It All Together}
        \begin{lstlisting}[language=Java]
    public class Main {
        public static void main(String[] args) {
            WeatherStation weatherStation = new WeatherStation();
            WeatherDisplay display = new WeatherDisplay();
            
            weatherStation.attach(display);
            
            // Simulate weather change
            weatherStation.setWeather("Sunny");
            weatherStation.setWeather("Rainy");
        }
    }
        \end{lstlisting}
    \end{block}
\end{frame}

\begin{frame}[fragile]
    \frametitle{Implementing the Observer Pattern - Conclusion}
    \begin{block}{Key Points}
        \begin{itemize}
            \item The Observer Pattern decouples the subject and observers, promoting a flexible and maintainable design.
            \item It is particularly useful in event-driven programming where multiple components need to react to changes.
            \item Supports the open/closed principle by allowing new observers to be added without modifying the subject.
        \end{itemize}
    \end{block}
    
    \begin{block}{Conclusion}
        The Observer Pattern enhances the modularity of code and simplifies the interaction between objects. 
        By following the outlined steps and understanding each component's role, you will be well-equipped to implement this pattern effectively.
    \end{block}
\end{frame}

\begin{frame}
    \frametitle{Comparing Singleton and Observer Patterns - Overview}
    \begin{block}{Overview of Software Design Patterns}
        Software design patterns are standard solutions to common problems in software design. 
        Two important patterns are \textbf{Singleton} and \textbf{Observer}, each serving different purposes.
    \end{block}
\end{frame}

\begin{frame}[fragile]
    \frametitle{Comparing Singleton and Observer Patterns - Singleton Pattern}
    \begin{block}{Singleton Pattern}
        \textbf{Definition:} The Singleton Pattern restricts a class to a single instance and provides a global point of access to it.
        
        \textbf{Use Cases:}
        \begin{itemize}
            \item Configuration Management: Ensures there is only one configuration object to manage application settings.
            \item Logging: Maintains a single log file for all parts of an application to allow consistent logging behavior.
            \item Connection Pooling: Controls access to a limited resource, such as a database connection.
        \end{itemize}
        
        \textbf{Example Implementation in Python:}
        \begin{lstlisting}[language=Python]
class Singleton:
    _instance = None

    def __new__(cls, *args, **kwargs):
        if not cls._instance:
            cls._instance = super(Singleton, cls).__new__(cls)
        return cls._instance
        \end{lstlisting}
    \end{block}
\end{frame}

\begin{frame}[fragile]
    \frametitle{Comparing Singleton and Observer Patterns - Observer Pattern}
    \begin{block}{Observer Pattern}
        \textbf{Definition:} The Observer Pattern defines a one-to-many dependency between objects, so when one object (the subject) changes state, 
        all its dependents (observers) are notified and updated automatically.
        
        \textbf{Use Cases:}
        \begin{itemize}
            \item Event Handling: Used in GUI frameworks where user actions (like clicks) trigger multiple updates.
            \item Subscription Systems: Enables real-time updates to subscribers when a publisher modifies data (e.g., stock prices).
            \item Dependency Injection: Notifies dependent components (observers) when a specific service or resource is available.
        \end{itemize}
        
        \textbf{Example Implementation in Python:}
        \begin{lstlisting}[language=Python]
class Subject:
    def __init__(self):
        self._observers = []

    def attach(self, observer):
        self._observers.append(observer)

    def notify(self):
        for observer in self._observers:
            observer.update()

class Observer:
    def update(self):
        print("Observer has been notified.")
        \end{lstlisting}
    \end{block}
\end{frame}

\begin{frame}
    \frametitle{Comparing Singleton and Observer Patterns - Key Differences}
    \begin{block}{Key Differences}
        \begin{center}
        \begin{tabular}{|l|l|l|}
            \hline
            \textbf{Aspect} & \textbf{Singleton Pattern} & \textbf{Observer Pattern} \\
            \hline
            Purpose & Ensure single instance of a class & Notify multiple observers about changes \\
            \hline
            Use Case & Centralized resource management & Dynamic updates in response to events \\
            \hline
            Instance Creation & Controlled through private constructor & Multiple observer instances can exist \\
            \hline
            Coupling & Low coupling to ensure accessibility & Higher coupling due to direct dependencies \\
            \hline
        \end{tabular}
        \end{center}
    \end{block}
\end{frame}

\begin{frame}
    \frametitle{Comparing Singleton and Observer Patterns - Choosing Patterns}
    \begin{block}{When to Favor One Over the Other}
        \begin{itemize}
            \item \textbf{Choose Singleton} when you need to control access to shared resources or require a single point of control.
            \item \textbf{Choose Observer} when you need to implement a publish/subscribe mechanism where changes in state need to be communicated 
            to multiple components effectively.
        \end{itemize}
    \end{block}
\end{frame}

\begin{frame}
    \frametitle{Comparing Singleton and Observer Patterns - Conclusion}
    \begin{block}{Conclusion}
        Understanding the differences and use cases for the Singleton and Observer patterns is crucial for effective software design 
        and ensuring appropriate patterns are applied based on application requirements.
    \end{block}
\end{frame}

\begin{frame}
    \frametitle{Best Practices in Using Design Patterns - Overview}
    \begin{block}{Why This Matters}
        Using design patterns wisely can significantly enhance the maintainability and flexibility of your software projects. However, misapplications can lead to complex designs and anti-patterns.
    \end{block}
\end{frame}

\begin{frame}
    \frametitle{Best Practices for Employing Design Patterns}
    \begin{enumerate}
        \item \textbf{Understand the Problem First}
            \begin{itemize}
                \item Design patterns are not one-size-fits-all solutions; they must fit the specific context of your project.
            \end{itemize}
            
        \item \textbf{Choose the Right Pattern}
            \begin{itemize}
                \item Familiarize yourself with various design patterns and their purposes.
                \item \textit{Example:} Use the Observer pattern in event-driven systems.
            \end{itemize}
    \end{enumerate}
\end{frame}

\begin{frame}
    \frametitle{Best Practices Continued}
    \begin{enumerate}
        \setcounter{enumi}{2}
        \item \textbf{Avoid Overusing Patterns}
            \begin{itemize}
                \item Applying too many design patterns can complicate designs and bloat code.
            \end{itemize}
        
        \item \textbf{Keep It Simple}
            \begin{itemize}
                \item Favor straightforward solutions over complex implementations.
            \end{itemize}
        
        \item \textbf{Document Your Use of Patterns}
            \begin{itemize}
                \item Use comments or markdown to explain the purpose and usage.
            \end{itemize}
    \end{enumerate}
\end{frame}

\begin{frame}
    \frametitle{Common Pitfalls}
    \begin{itemize}
        \item \textbf{Misapplication:} Using a pattern without understanding can lead to suboptimal outcomes.
        \item \textbf{Anti-Patterns:} Overengineering can create harmful structures (e.g., God Object).
        \item \textbf{Ignoring Context:} A pattern may not work in another system if context is disregarded.
    \end{itemize}
\end{frame}

\begin{frame}[fragile]
    \frametitle{Code Snippet Example: Singleton Pattern in Python}
    \begin{lstlisting}[language=Python]
class Singleton:
    _instance = None

    def __new__(cls, *args, **kwargs):
        if not cls._instance:
            cls._instance = super(Singleton, cls).__new__(cls)
        return cls._instance

# Usage
singleton1 = Singleton()
singleton2 = Singleton()

assert singleton1 is singleton2  # Both references point to the same instance
    \end{lstlisting}
\end{frame}

\begin{frame}
    \frametitle{Conclusion and Key Takeaways}
    \begin{itemize}
        \item Use design patterns thoughtfully to enhance project quality.
        \item Key Takeaways:
            \begin{itemize}
                \item Know your problem context.
                \item Match patterns appropriately.
                \item Strive for simplicity in design.
                \item Document and review pattern usage.
            \end{itemize}
    \end{itemize}
\end{frame}

\begin{frame}[fragile]
    \frametitle{Conclusion and Further Resources - Recap of Key Concepts}

    \begin{block}{What Are Software Design Patterns?}
        Software design patterns are reusable solutions to common problems that occur within a given context in software design. 
        They provide a template for how to solve problems and enable developers to communicate more effectively.
    \end{block}

    \begin{block}{Types of Design Patterns}
        \begin{itemize}
            \item \textbf{Creational Patterns:} Deal with object creation mechanisms.
            \begin{itemize}
                \item \textbf{Singleton Pattern:} Ensures a class has only one instance and provides a global point of access.
                \item \textbf{Factory Method Pattern:} Defines an interface for creating an object but lets subclasses alter the type of objects that will be created.
            \end{itemize}

            \item \textbf{Structural Patterns:} Concerned with how classes and objects are composed to form larger structures.
            \begin{itemize}
                \item \textbf{Adapter Pattern:} Allows incompatible interfaces to work together.
                \item \textbf{Facade Pattern:} Provides a simplified interface to a complex subsystem.
            \end{itemize}

            \item \textbf{Behavioral Patterns:} Focus on communication between objects.
            \begin{itemize}
                \item \textbf{Observer Pattern:} Defines a one-to-many dependency between objects so that when one object changes state, all its dependents are notified.
                \item \textbf{Strategy Pattern:} Enables selecting an algorithm's behavior at runtime.
            \end{itemize}
        \end{itemize}
    \end{block}
\end{frame}

\begin{frame}[fragile]
    \frametitle{Conclusion and Further Resources - Best Practices and Key Takeaways}

    \begin{block}{Best Practices}
        \begin{itemize}
            \item Understand the context before applying a design pattern.
            \item Avoid unnecessary complexity — Choose patterns that fit the simplicity of the problem.
            \item Refactor code as needed, in conjunction with design patterns, to improve maintainability.
        \end{itemize}
    \end{block}

    \begin{block}{Key Takeaways}
        \begin{itemize}
            \item Design patterns are essential tools that enhance code reusability and flexibility.
            \item Choose the right pattern to fit the specific problem context; design patterns are not a one-size-fits-all solution.
            \item Continuous learning and practice with design patterns will improve software development skills.
        \end{itemize}
    \end{block}
\end{frame}

\begin{frame}[fragile]
    \frametitle{Conclusion and Further Resources - Further Resources for Learning}

    \begin{block}{Books}
        \begin{itemize}
            \item \textit{Design Patterns: Elements of Reusable Object-Oriented Software} by Erich Gamma, Richard Helm, Ralph Johnson, and John Vlissides (Gang of Four).
            \item \textit{Head First Design Patterns} by Eric Freeman and Bert Bates - An engaging, visually rich approach to design patterns.
        \end{itemize}
    \end{block}

    \begin{block}{Online Platforms}
        \begin{itemize}
            \item \textbf{Coursera:} Courses on software engineering best practices including design patterns.
            \item \textbf{Pluralsight:} Video tutorials specifically addressing modern design patterns in various programming languages.
        \end{itemize}
    \end{block}

    \begin{block}{Communities and Documentation}
        \begin{itemize}
            \item \textbf{Stack Overflow:} Participate in discussions and get answers to specific design pattern questions.
            \item \textbf{GitHub:} Explore repositories that implement design patterns in various programming languages.
            \item \textbf{Refactoring Guru:} Offers explanations and visual examples of different design patterns.
            \item \textbf{Vega's Patterns Catalog:} A concise overview of design patterns with examples and use cases.
        \end{itemize}
    \end{block}

    \begin{block}{Closing Note}
        Continually explore and practice with design patterns to ensure you become proficient at incorporating them into your software development projects. Happy coding!
    \end{block}
\end{frame}


\end{document}